\chapter{Dataset and Ground Truth}
\label{ch:dataset}

\section{Data Sources}

The dataset used in this study was assembled from two sources:

\begin{enumerate}
  \item \textbf{Real production-stage PDFs}: IKEA of Sweden provided a set of label PDF
  files in various states of completion—approved templates, draft revisions, and known-faulty
  samples flagged by the internal \ac{QA} team. These files represent the actual artefacts
  that the system is intended to process in production. Due to confidentiality requirements,
  all product names, article numbers, and supplier information in published materials have
  been anonymised.

  \item \textbf{Synthetically modified labels}: To increase the volume and diversity of
  non-compliant samples, a set of compliant templates was programmatically modified using
  PyMuPDF to introduce controlled violations. Violation types included: (a) removing
  required fields, (b) corrupting the article number format, (c) replacing the country-of-origin
  string, (d) removing the DataMatrix barcode, and (e) corrupting GS1 AI fields in the
  DataMatrix payload.
\end{enumerate}

\section{Dataset Composition}

The final dataset consists of 238 annotated label images drawn from 28 of the 34 label
types in the schema corpus (the remaining 6 types could not be represented due to
unavailability of real samples and their complexity making reliable synthetic generation
infeasible within the project timeline).

\begin{table}[h]
\centering
\caption{Dataset split and class distribution.}
\label{tab:dataset_split}
\begin{tabular}{lcccc}
\toprule
\textbf{Split} & \textbf{Total} & \textbf{Compliant} & \textbf{Non-compliant} & \textbf{\%} \\
\midrule
Train       & 166 & 91  & 75  & 69.7\% \\
Validation  & 36  & 20  & 16  & 15.1\% \\
Test        & 36  & 20  & 16  & 15.1\% \\
\midrule
\textbf{Total} & \textbf{238} & \textbf{131} & \textbf{107} & 100\% \\
\bottomrule
\end{tabular}
\end{table}

The 55\%/45\% class split is intentional: non-compliant labels are slightly less frequent in
the real production pipeline, but over-representing non-compliant samples in the training
data ensures the validator is exposed to diverse failure modes.

\Cref{tab:label_type_distribution} shows the five most represented label types in the dataset
alongside their rule corpus statistics.

\begin{table}[h]
\centering
\caption{Five most represented label types and their rule complexity.}
\label{tab:label_type_distribution}
\begin{tabular}{lcccc}
\toprule
\textbf{Label Type} & \textbf{Samples} & \textbf{Required fields} & \textbf{Optional fields} & \textbf{Rules} \\
\midrule
1R35-IDADM\_IDBDM & 24 & 18 & 7  & 11 \\
1R10-IDADM\_IDBDM & 22 & 18 & 6  & 10 \\
1C70-CAA          & 18 & 16 & 3  & 4  \\
L10555-MIADM      & 16 & 16 & 2  & 5  \\
1C50-IDBDM        & 14 & 15 & 3  & 2  \\
\midrule
Other (23 types)  & 144 & --- & --- & --- \\
\bottomrule
\end{tabular}
\end{table}

\section{Ground-Truth Rule Schemas}

The 34 JSON schema files form the ground truth of the validation system. Each file was
authored based on IKEA of Sweden's internal label design specification document (a
90-page PDF covering all label types, barcode specifications, zone definitions, and
variant rules). The process of converting the specification into machine-readable schemas
involved:

\begin{enumerate}
  \item Reading each label type description and extracting the list of required and optional
  fields.
  \item Translating layout rules (e.g.\ ``the DataMatrix must be in the centre zone'') into
  JSON rule objects.
  \item Encoding the GS1 AI requirements for each barcode type as \texttt{ai\_field\_present}
  rules.
  \item Encoding content constraints (article number format, origin country whitelist) as
  \texttt{regex\_match} or \texttt{exact\_match} rules.
\end{enumerate}

Across all 34 schemas, the corpus contains:
\begin{itemize}
  \item \textbf{342 required-field declarations} (mean: 10.1 per schema, range: 2--18).
  \item \textbf{140 optional-field declarations} (mean: 4.1 per schema, range: 0--7).
  \item \textbf{107 validation rules} (mean: 3.1 per schema, range: 2--11).
\end{itemize}

The schema design was reviewed by two domain experts at IKEA of Sweden and one revision
cycle was performed to correct ambiguities and add rules that were implicit in the
specification but not initially captured.

\section{Annotation Methodology}

Each label in the dataset was annotated at the field level using a semi-automated process:

\begin{enumerate}
  \item The label was processed by the pipeline in its current state.
  \item The extracted fields were displayed in the Streamlit \ac{UI}.
  \item A human annotator reviewed and corrected any extraction errors using the
  built-in correction interface.
  \item The corrected extraction was exported as an NDJSON training record via
  \texttt{/api/export-jsonl}.
  \item Compliance ground truth (valid/invalid) was set by the person who created the
  label (for synthetic samples) or confirmed by the IKEA \ac{QA} team (for real samples).
\end{enumerate}

Inter-annotator agreement was measured on a random sample of 40 labels annotated
independently by two authors. Cohen's kappa for the overall compliance label was~0.84,
indicating strong agreement. Field-level extraction disagreements were most frequent on the
\texttt{dimensions\_metric} field, where the format varies across label types.

\section{DPI and Rendering Quality}

All PDF labels were rendered at 300~DPI for consistency. Tests at 150~DPI and 200~DPI
showed measurable drops in barcode decoding rates (ITF-14 decode rate fell from~94.4\%
at 300~DPI to~81.2\% at 150~DPI) and OCR character accuracy (from~89.1\% to~81.6\%).
Raster input images were accepted at their native resolution; images below~200~DPI were
flagged with a quality warning in the validation report.

\section{Privacy and Data Handling}

All label data used in this study was provided under a data-sharing agreement with IKEA of
Sweden. The following measures were applied:
\begin{itemize}
  \item No labels containing personally identifiable information were included.
  \item Article numbers and internal item codes were masked with a one-way hash in any
  exported dataset artefacts.
  \item The dataset is not publicly released; access is available to Linnaeus University
  researchers via a signed data-sharing agreement with IKEA of Sweden.
\end{itemize}
